\chapter{Methodology}
\label{sec:chapterMethodology}

This chapter describes the experimental methodology adopted in the thesis. Chapter~\ref{sec:chapterTheoreticalBackground} (p.~\pageref{sec:chapterTheoreticalBackground}) defines the learning algorithms, imbalance interventions, evaluation metrics, threshold rules, hyperparameter optimisation procedure, and feature-attribution estimators. The present chapter explains how those elements were configured and applied in this study.

Section~\ref{sec:objective} links the research questions to the overall experimental design. Sections~\ref{sec:datasets}--\ref{sec:models} introduce the datasets and models, while Sections~\ref{sec:imbalance_strategies}--\ref{sec:protocol} describe the imbalance interventions and leakage-free evaluation protocol. Sections~\ref{sec:threshold_study}--\ref{sec:tuning} specify threshold selection and hyperparameter optimisation. Section~\ref{sec:cross_domain} defines the BAF transfer experiment, and Sections~\ref{sec:metrics}--\ref{sec:interpretability} document the reported metrics, reproducibility provisions, and interpretability analyses.

% ──────────────────────────────────────────────────────────────────────────
\section{Research Objective and Experimental Design}
\label{sec:objective}
% ──────────────────────────────────────────────────────────────────────────

This thesis investigates fraud detection in financial transactions under extreme class imbalance. The work is structured around six complementary research contributions:

\begin{enumerate}
  \item \textbf{Unified model comparison.} Compare classical machine learning models and a transformer-based tabular model across distinct fraud datasets using a shared leakage-free outer holdout, common metrics, and model-appropriate validation procedures.

  \item \textbf{Threshold selection as a first-class experimental variable.} Examine how the decision-threshold rule affects model rankings and operational metrics using validation-only selection.

  \item \textbf{Full factorial analysis of imbalance handling.} Evaluate seven imbalance strategies $\times$ five supervised model families $\times$ two datasets in a structured design, documenting how each intervention changes the observed metrics.

  \item \textbf{Cross-domain generalisation under controlled BAF conditions.} Assess robustness by training on BAF Base and evaluating, without retraining, on the five BAF Variants that alter group-size, prevalence, separability, or temporal conditions.

  \item \textbf{SHAP feature-set consistency.} Compare top-feature sets from architecture-compatible SHAP explainers for Logistic Regression, LightGBM, and CatBoost, and assess the stability of the Base-trained LGBM top-feature set across BAF Variant test distributions.

  \item \textbf{FT-Transformer attention diagnostics.} Characterise the concentration of last-layer, head-averaged \texttt{[CLS]} attention across BAF feature tokens as a descriptive diagnostic that is separate from SHAP feature attribution.
\end{enumerate}

Contributions~5 and~6 form the two complementary components of RQ5. Table~\ref{tab:rq_design_map} shows how all five research questions are operationalised and where their evidence is reported.

\begin{table}[htbp]
\centering
\small
\caption{Alignment between research questions, experimental operations, and Results sections.}
\label{tab:rq_design_map}
\begin{tabular}{p{0.08\linewidth} p{0.52\linewidth} p{0.27\linewidth}}
\toprule
\textbf{RQ} & \textbf{Experimental operationalisation} & \textbf{Primary evidence} \\
\midrule
RQ1 & Baseline comparison of the six model families under the shared outer holdout and model-specific inner validation. & Section~\ref{sec:baseline_results} \\
RQ2 & Comparison of fixed, max-$F_1$, max-$F_2$, and precision-constrained threshold rules selected within DEV. & Section~\ref{sec:threshold_sensitivity} \\
RQ3 & Factorial comparison of seven imbalance strategies across five supervised models and two datasets, followed by a focused FT-Transformer--CatBoost comparison. & Sections~\ref{sec:factorial_results} and~\ref{sec:transformer_imbalance} \\
RQ4 & Evaluation of fixed BAF Base-trained models and thresholds on Variant test distributions without retraining. & Section~\ref{sec:cross_domain_results} \\
RQ5 & SHAP top-feature-set comparisons for compatible explainers and a separate descriptive FT-Transformer attention diagnostic. & Sections~\ref{sec:shap_results} and~\ref{sec:attention_diagnostics} \\
\bottomrule
\end{tabular}
\end{table}

Figure~\ref{fig:pipeline} provides a high-level overview of the end-to-end experimental pipeline described in the remainder of this chapter.

% ═══════════════════════════════════════════════════════════════════════════
%  FIGURE 4.1 — End-to-End Experimental Pipeline
% ═══════════════════════════════════════════════════════════════════════════

\definecolor{clrData}{HTML}{4472C4}      % Steel blue
\definecolor{clrSplit}{HTML}{5B9BD5}     % Light blue
\definecolor{clrDev}{HTML}{FFF2CC}       % Light yellow (DEV zone)
\definecolor{clrDevBorder}{HTML}{BF8F00} % Dark gold (DEV border)
\definecolor{clrPrep}{HTML}{ED7D31}      % Orange
\definecolor{clrTune}{HTML}{FFC000}      % Gold
\definecolor{clrThresh}{HTML}{F4B183}    % Light orange
\definecolor{clrFact}{HTML}{E2EFDA}      % Light green (factorial zone)
\definecolor{clrFactBorder}{HTML}{548235}% Dark green
\definecolor{clrTest}{HTML}{D6E4F0}      % Light blue (TEST zone)
\definecolor{clrTestBorder}{HTML}{2E75B6}% Dark blue
\definecolor{clrFinal}{HTML}{2E75B6}     % Blue
\definecolor{clrEval}{HTML}{00B050}      % Green
\definecolor{clrShap}{HTML}{7030A0}      % Purple

\par\medskip\noindent\begin{minipage}{\linewidth}\centering
\begin{tikzpicture}[
    node distance=0.6cm,
    >={Stealth[length=3mm]},
    box/.style={
        rectangle, rounded corners=3pt, minimum width=10.5cm, minimum height=1.1cm,
        text=white, font=\footnotesize\bfseries, align=center,
        drop shadow={shadow xshift=1pt, shadow yshift=-1pt}
    },
    smallbox/.style={
        rectangle, rounded corners=3pt, minimum width=8cm, minimum height=0.9cm,
        font=\footnotesize\bfseries, align=center,
        drop shadow={shadow xshift=1pt, shadow yshift=-1pt}
    },
    arrow/.style={->, thick, black!70},
    zone/.style={rectangle, rounded corners=5pt, draw=#1, line width=1.2pt, dashed, inner sep=10pt},
]

% ── Datasets ──
\node[box, fill=clrData] (data)
    {1.\ Datasets\\[-2pt]
     \fontsize{7}{8}\selectfont ULB 2013 ($N{=}284\text{K}$, $0.17\%$ fraud)
     \;\;$\cdot$\;\; BAF Suite ($N{=}1\text{M}$, ${\approx}1.1\%$ fraud)};

% ── Split ──
\node[box, fill=clrSplit, below=of data] (split)
    {2.\ Stratified Holdout Split\\[-2pt]
     \fontsize{7}{8}\selectfont 80\% DEV \;/\; 20\% TEST \;\;$\cdot$\;\; stratified by target};

% ── DEV zone ──
\node[box, fill=clrPrep, below=0.85cm of split] (prep)
    {3.\ Preprocessing\\[-2pt]
     \fontsize{7}{8}\selectfont Classical/OCSVM: impute, scale, one-hot \;\;$\cdot$\;\; FT: scale numerical, ordinal categories};

\node[box, fill=clrTune, below=of prep] (tune)
    {4.\ Supervised-model Hyperparameter Tuning (Optuna)\\[-2pt]
     \fontsize{7}{8}\selectfont TPE sampler \;\;$\cdot$\;\; 50 trials \;\;$\cdot$\;\; PR-AUC objective};

\node[box, fill=clrThresh, text=black, below=of tune] (thresh)
    {5.\ Validation and Threshold Selection\\[-2pt]
     \fontsize{7}{8}\selectfont Classical/OCSVM: 5-fold median $\tau$ \;\;$\cdot$\;\; FT: single DEV holdout $\tau$};

% ── Factorial annotation ──
\node[smallbox, fill=clrFact!80, text=clrFactBorder, draw=clrFactBorder,
      line width=0.8pt, below=of thresh] (fact)
    {$7$ strategies $\times$ $5$ models $= 35$ configurations per dataset};

% DEV zone background
\begin{scope}[on background layer]
    \node[zone=clrDevBorder, fill=clrDev, fit=(prep)(tune)(thresh)(fact),
          label={[font=\footnotesize\bfseries, text=clrDevBorder]above left:DEV}] (devzone) {};
\end{scope}

% ── TEST zone ──
\node[box, fill=clrFinal, below=1.0cm of fact] (final)
    {6.\ Final Training\\[-2pt]
     \fontsize{7}{8}\selectfont Full training partition \;\;$\cdot$\;\; selected configuration};

\node[box, fill=clrEval, below=of final] (eval)
    {7.\ Held-Out Test Evaluation\\[-2pt]
     \fontsize{7}{8}\selectfont Apply validation-derived $\tau$ \;\;$\cdot$\;\; Bootstrap 95\% CI ($B{=}1000$)};

% TEST zone background
\begin{scope}[on background layer]
    \node[zone=clrTestBorder, fill=clrTest, fit=(final)(eval),
          label={[font=\footnotesize\bfseries, text=clrTestBorder]above left:TEST}] (testzone) {};
\end{scope}

% ── SHAP ──
\node[box, fill=clrShap, below=1.0cm of eval] (shap)
    {8.\ Interpretability Analyses\\[-2pt]
     \fontsize{7}{8}\selectfont SHAP feature sets and local cases \;\;$\cdot$\;\; FT attention diagnostic};

% ── Arrows ──
\draw[arrow] (data) -- (split);
\draw[arrow] (split) -- (prep);
\draw[arrow] (prep) -- (tune);
\draw[arrow] (tune) -- (thresh);
\draw[arrow] (thresh) -- (fact);
\draw[arrow] (fact) -- (final);
\draw[arrow] (final) -- (eval);
\draw[arrow] (eval) -- (shap);

\end{tikzpicture}
\captionof{figure}[End-to-end experimental pipeline for fraud detection]{End-to-end experimental pipeline for fraud detection under extreme class imbalance. Stages~3--5 operate exclusively on the DEV partition (80\%): classical models and OCSVM use five-fold validation, whereas FT-Transformer uses a single internal holdout. Stages~6--7 provide one evaluation on the untouched TEST partition (20\%). Stage~8 comprises SHAP for the linear and tree-based models and a separate attention diagnostic for FT-Transformer.}
\label{fig:pipeline}
\end{minipage}\par\medskip
\clearpage

% ──────────────────────────────────────────────────────────────────────────
\section{Datasets}
\label{sec:datasets}
% ──────────────────────────────────────────────────────────────────────────

\paragraph{Credit Card Fraud Detection 2013 (ULB).}
The ULB dataset contains European credit card transactions with anonymised features (PCA components V1--V28, plus \texttt{Time} and \texttt{Amount}) and a binary target indicating fraud. The dataset exhibits extreme class imbalance ($\approx$0.17\% fraud).

\paragraph{Bank Account Fraud (BAF) Suite (NeurIPS 2022).}
The BAF suite contains a Base dataset and five variants (Variant~I--V) designed to stress-test models under controlled bias and distributional conditions. Each dataset comprises $\approx$1\,000\,000 instances with mixed numerical and categorical variables (5 categorical features: \texttt{payment\_type}, \texttt{employment\_status}, \texttt{housing\_status}, \texttt{source}, \texttt{device\_os}) and a binary label \texttt{fraud\_bool} ($\approx$1.1\% fraud in Base). The raw Base file contains 32 columns, including the target and the \texttt{month} field. The pipeline excludes both fields from the predictors, yielding 30 modelling features. Variants~III and~V additionally contain \texttt{x1} and \texttt{x2}; these synthetic columns are also excluded so that all Base-to-Variant evaluations use the same 30-feature input.

\paragraph{Cross-dataset compatibility note.}
The ULB and BAF datasets have incompatible feature spaces (PCA-anonymised vs.\ semantic features). Therefore, cross-dataset transfer between ULB and BAF is not feasible. Cross-domain generalisation is studied \emph{within} the BAF suite, where a shared modelling feature space enables Base-trained models to be evaluated on each controlled variant.

% ──────────────────────────────────────────────────────────────────────────
\section{Models}
\label{sec:models}
% ──────────────────────────────────────────────────────────────────────────

Six model families are evaluated: five supervised classifiers --- Logistic Regression (LR), Random Forest (RF), LightGBM (LGBM)~\cite{ke2017lightgbm}, CatBoost (CB)~\cite{prokhorenkova2018catboost}, and the Feature Tokenizer Transformer (FT-Transformer)~\cite{gorishniy2021revisiting} --- together with One-Class SVM (OCSVM) as an unsupervised anomaly-detection baseline. Definitions of each architecture are given in Section~\ref{sec:tb:models}.

The selection covers four supervised approaches --- linear classification (LR), bagged trees (RF), gradient-boosted trees (LGBM and CB), and an attention-based deep network (FT-Transformer) --- together with the OCSVM reference. The two gradient-boosted implementations are retained separately: LightGBM provides a computationally efficient tabular model, whereas CatBoost uses a distinct boosting procedure. OCSVM quantifies performance without supervised fraud labels and is therefore excluded from the imbalance-strategy and class-weighting comparisons. Including FT-Transformer alongside the four classical supervised families allows their observed responses to the implemented imbalance interventions to be compared without assuming that any difference is inherent to the architecture.

\paragraph{Model-specific preprocessing.}
The classical models and OCSVM use a shared \texttt{ColumnTransformer}: numerical variables are mean-imputed and standardised, while categorical variables are mode-imputed and one-hot encoded. Consequently, although CatBoost supports native categorical processing in general, that capability is not used in this experiment. FT-Transformer follows a separate input path in which numerical variables are mean-imputed and standardised and categorical variables are ordinally encoded as integer embedding indices. Each preprocessor is fitted only on the relevant training partition.

% ──────────────────────────────────────────────────────────────────────────
\section{Imbalance Handling Strategies}
\label{sec:imbalance_strategies}
% ──────────────────────────────────────────────────────────────────────────

To address extreme class imbalance, seven strategies are evaluated for the five supervised model families (LR, RF, LGBM, CB, and FT-Transformer): an unmodified baseline (\textbf{None}); random undersampling (\textbf{RUS}) and random oversampling (\textbf{ROS}); \textbf{SMOTE}~\cite{chawla2002smote} and its two cleaning-augmented variants, \textbf{SMOTE+Tomek Links} and \textbf{SMOTEENN}~\cite{batista2004balancing}; and cost-sensitive \textbf{Class Weights}. The selection covers undersampling, oversampling, and cost-sensitive adjustment~\cite{hegarcia2009learning, krawczyk2016learning}, together with the unmodified reference needed to measure changes associated with each intervention. Section~\ref{sec:tb:imbalance} defines the algorithms.

OCSVM is evaluated separately (no resampling or class weighting), given its one-class anomaly-detection paradigm. The full factorial design yields $7 \text{ strategies} \times 5 \text{ supervised models} = 35$ configurations per dataset (plus OCSVM), enabling a structured analysis of the \emph{strategy~$\times$~model~$\times$~dataset} interaction.

\paragraph{Resampling representation.}
The implemented samplers operate after model-specific preprocessing. For the classical models, SMOTE therefore interpolates the standardised numerical variables and the one-hot-encoded categorical columns. For FT-Transformer, the sampler receives concatenated standardised numerical values and ordinal category codes; synthetic categorical values are subsequently converted to integer embedding indices. Neither path uses a categorical-aware sampler such as SMOTENC. The resulting SMOTE comparisons describe the behaviour of this implemented pipeline and cannot isolate an architecture-only effect from the representation used during synthetic resampling.

% ──────────────────────────────────────────────────────────────────────────
\section{Leakage-Free Evaluation Protocol}
\label{sec:protocol}
% ──────────────────────────────────────────────────────────────────────────

A central methodological concern in fraud detection is preventing data leakage: information from the test set must never influence model selection, preprocessing, resampling, or threshold calibration. This section formalises the leakage-free protocol adopted throughout this thesis. Figure~\ref{fig:cv_protocol} illustrates the inner workings of a single cross-validation fold.

\subsection{Outer Split (Holdout Test Set)}
For each dataset, an outer stratified split creates an 80\% development partition (DEV) and a 20\% holdout test partition (TEST). TEST is not used to fit preprocessing, resample data, select hyperparameters, or select decision thresholds. Post-hoc analyses that use TEST are conducted only after the predictive evaluation has been fixed.

\subsection{Inner Validation (Model Selection and Thresholding)}
For the classical supervised models and OCSVM, selection occurs exclusively within the DEV partition via stratified $k$-fold cross-validation ($k=5$). For each fold:
\begin{enumerate}
  \item Fit the preprocessor on the fold's training subset only.
  \item Apply the imbalance strategy only on the fold's training subset.
  \item Fit the model.
  \item Evaluate on the fold's validation subset: compute PR-AUC and obtain predicted scores.
  \item Select a decision threshold $\tau$ by maximising the chosen criterion (see Section~\ref{sec:threshold_study}) on the validation subset.
  \item Using $\tau$, compute $F_1$, $F_2$, and confusion matrix counts $\{TN, FP, FN, TP\}$.
\end{enumerate}

For OCSVM, the model is fitted only on legitimate examples from each training fold, while the corresponding validation fold retains both classes for threshold selection and evaluation. FT-Transformer replaces five-fold cross-validation with one stratified 80/20 holdout inside DEV. Its preprocessing, early stopping, hyperparameter selection, and $F_2$ threshold selection use only this internal training/validation split. This model-specific difference reduces computational cost but means that fold dispersion and median thresholds are available only for the classical models and OCSVM.

\subsection{Final Training and Test Evaluation}
For each model--strategy--dataset configuration, the final model is fitted on the full DEV partition. For the classical models and OCSVM, the decision threshold $\tau$ applied to TEST is the median of the five validation-fold thresholds. For FT-Transformer, it is the threshold selected on the single internal validation split. In both paths, TEST remains untouched until the final predictive evaluation.

% ═══════════════════════════════════════════════════════════════════════════
%  FIGURE 4.2 — Leakage-Free Inner Cross-Validation Protocol
% ═══════════════════════════════════════════════════════════════════════════

\definecolor{clrTrainFold}{HTML}{DAE8FC}    % Light blue (train fold)
\definecolor{clrValFold}{HTML}{D5E8D4}      % Light green (validation fold)
\definecolor{clrBarrier}{HTML}{CC0000}       % Red (leakage barrier)
\definecolor{clrStep}{HTML}{4472C4}         % Blue (process steps)
\definecolor{clrValStep}{HTML}{548235}      % Green (validation steps)
\definecolor{clrWarn}{HTML}{CC0000}         % Red (warning annotations)

\par\medskip\noindent\begin{minipage}{\linewidth}\centering
\begin{tikzpicture}[
    >={Stealth[length=2.5mm]},
    node distance=0.5cm,
    procstep/.style={
        rectangle, rounded corners=2pt, minimum width=3.8cm, minimum height=0.8cm,
        font=\footnotesize\bfseries, align=center, text=white,
        drop shadow={shadow xshift=0.8pt, shadow yshift=-0.8pt}
    },
    trainstep/.style={procstep, fill=clrStep},
    valstep/.style={procstep, fill=clrValStep},
    annot/.style={font=\scriptsize, text=black!60, align=center},
    warnnote/.style={font=\scriptsize\itshape, text=clrWarn, align=center},
    arrow/.style={->, thick, black!60},
    zone/.style={rectangle, rounded corners=4pt, draw=#1, line width=1.2pt, inner sep=10pt},
]

% ── Title ──
\node[font=\small\bfseries, text=black!80] (title) {One Classical/OCSVM Inner CV Fold ($k$ of $K{=}5$)};

% ── TRAIN FOLD (left side) ──
\node[trainstep, below=0.9cm of title, xshift=-5.2cm] (fitpre)
    {1.\ Fit Preprocessor};
\node[annot, below=0.15cm of fitpre] {on train subset only};

\node[trainstep, below=1.0cm of fitpre] (resample)
    {2.\ Apply Resampling};
\node[annot, below=0.15cm of resample] {RUS / ROS / SMOTE / \ldots};

\node[trainstep, below=1.0cm of resample] (train)
    {3.\ Train Model};
\node[annot, below=0.15cm of train] {on resampled train data};

% Train zone background
\begin{scope}[on background layer]
    \node[zone=clrStep!70, fill=clrTrainFold, fit=(fitpre)(resample)(train),
          inner sep=16pt,
          label={[font=\footnotesize\bfseries, text=clrStep]above:TRAIN FOLD}] (trainzone) {};
\end{scope}

% ── VALIDATION FOLD (right side) ──
\node[valstep, below=0.9cm of title, xshift=5.2cm] (transform)
    {4.\ Transform Only};
\node[annot, below=0.15cm of transform] {preprocessor fitted on train};

\node[valstep, below=1.0cm of transform] (score)
    {5.\ Score Predictions};
\node[annot, below=0.15cm of score] {PR-AUC \;\;$\cdot$\;\; predicted scores};

\node[valstep, below=1.0cm of score] (threshold)
    {6.\ Select Threshold $\tau$};
\node[annot, below=0.15cm of threshold] {maximise $F_2$ on validation};

% Validation zone background
\begin{scope}[on background layer]
    \node[zone=clrValStep!70, fill=clrValFold, fit=(transform)(score)(threshold),
          inner sep=16pt,
          label={[font=\footnotesize\bfseries, text=clrValStep]above:VALIDATION FOLD}] (valzone) {};
\end{scope}

% ── LEAKAGE BARRIER (centred between zones) ──
\coordinate (barTop) at ($(trainzone.north east)!0.5!(valzone.north west)+(0,0.4)$);
\coordinate (barBot) at ($(trainzone.south east)!0.5!(valzone.south west)+(0,-0.4)$);
\draw[clrBarrier, line width=2.5pt, dashed] (barTop) -- (barBot);
\node[font=\footnotesize\bfseries, text=clrBarrier, rotate=90, anchor=south]
    at ($(barTop)!0.5!(barBot)+(0.35,0)$) {LEAKAGE BARRIER};

% ── Arrows within train ──
\draw[arrow] (fitpre) -- (resample);
\draw[arrow] (resample) -- (train);

% ── Arrows within validation ──
\draw[arrow] (transform) -- (score);
\draw[arrow] (score) -- (threshold);

% ── Cross-zone arrows (train → validation) ──
\draw[->, thick, clrStep!60, densely dotted]
    (fitpre.east) -- ++(0.55,0) |- (transform.west)
    node[pos=0.25, above, font=\scriptsize, text=clrStep] {fitted preprocessor};
\draw[->, thick, clrStep!60, densely dotted]
    (train.east) -- ++(0.55,0) |- (score.west)
    node[pos=0.25, below, font=\scriptsize, text=clrStep] {trained model};

% ── Warning annotations ──
\node[warnnote, below=0.5cm of trainzone.south] (warn1)
    {$\triangleright$ No fit on validation data};
\node[warnnote, below=0.5cm of valzone.south] (warn2)
    {$\triangleright$ No resampling on validation data};

% ── Output ──
\node[rectangle, rounded corners=3pt, fill=black!8, draw=black!30,
      minimum width=8cm, minimum height=0.7cm,
      font=\footnotesize\bfseries, align=center,
      below=1.2cm of warn1.south east] (output)
    {Output:\; PR-AUC$_k$ \quad $\tau_k$ \quad $F_{1,k}$ \quad $F_{2,k}$ \quad $\{TN,FP,FN,TP\}_k$};

\draw[arrow] (threshold.south) |- (output.east);

% ── Final aggregation note ──
\node[annot, below=0.2cm of output]
    {Repeat for $k = 1, \ldots, K$. \quad Final $\tau = \mathrm{median}(\tau_1, \ldots, \tau_K)$.};

\end{tikzpicture}
\captionof{figure}[Leakage-free inner cross-validation protocol]{Leakage-free inner cross-validation protocol for the classical supervised models and OCSVM. Within each fold, preprocessing is fitted and any applicable resampling is performed exclusively on the training subset (left); the validation subset (right) is only transformed and scored. The dashed red line represents the leakage barrier, and the final decision threshold is the median across all $K{=}5$ folds. FT-Transformer instead uses the separate DEV holdout procedure described in the text.}
\label{fig:cv_protocol}
\end{minipage}\par\medskip

\clearpage

% ──────────────────────────────────────────────────────────────────────────
\section{Threshold Selection Study}
\label{sec:threshold_study}
% ──────────────────────────────────────────────────────────────────────────

Threshold selection is often implicit in fraud detection studies, and selecting a threshold on the test set constitutes data leakage. This thesis instead compares four explicit rules: the fixed default $\tau{=}0.5$; the validation-optimal $F_1$ threshold ($\tau^{*}=\arg\max_{\tau} F_1(\tau)$), which gives precision and recall equal weight; the validation-optimal $F_2$ threshold ($\tau^{*}=\arg\max_{\tau} F_2(\tau)$), which assigns greater weight to recall; and a precision-constrained rule that selects the smallest $\tau$ for which $\text{Precision}(\tau) \geq 0.5$ on validation. The last rule represents a minimum validation-precision target, not a guarantee that test precision will remain above 0.5 under sampling variation or distribution shift. Section~\ref{sec:tb:thresholds} defines these rules, and Section~\ref{sec:tb:metrics} defines $F_{\beta}$.

For OCSVM, these rules operate on the negated decision function rather than a calibrated probability. Its fixed value of 0.5 is retained as a numerical sensitivity reference but is not equivalent to a 50\% posterior-probability decision.

For the classical models and OCSVM, each rule is evaluated on the five validation folds and the median threshold is applied to the saved TEST scores. For FT-Transformer, thresholds are obtained from a separately reconstructed internal DEV split by training a validation model for the baseline-selected number of epochs; the resulting thresholds are then applied to the saved baseline TEST scores. TEST labels are never used to select a threshold. Because the baseline FT-Transformer threshold came from its original early-stopping validation run, its baseline max-$F_2$ operating point and the max-$F_2$ point reconstructed for this threshold study need not be numerically identical. The threshold study reports all four study-specific operating points to quantify changes in $F_1$, $F_2$, and workload metrics.

% ──────────────────────────────────────────────────────────────────────────
\section{Hyperparameter Tuning Strategy}
\label{sec:tuning}
% ──────────────────────────────────────────────────────────────────────────

Hyperparameter tuning is performed only for the baseline setting (\texttt{strategy=None}) to keep the factorial study computationally feasible and internally consistent. The baseline-selected hyperparameters are reused for every imbalance strategy within the same model and dataset. Model capacity is therefore held fixed while the planned imbalance intervention changes, although the comparison remains conditional on the baseline search result.

\subsection{Classical Models}

Tuning uses \textbf{Optuna}~\cite{akiba2019optuna} with the Tree-structured Parzen Estimator (TPE) sampler and a \texttt{MedianPruner} that terminates unpromising trials on the basis of per-fold intermediate results; Section~\ref{sec:tb:hpo} describes both methods. Each study runs 50~trials and optimises PR-AUC through five-fold stratified cross-validation. For LightGBM and CatBoost, early stopping with a patience of $50$ rounds limits unnecessary boosting iterations and controls overfitting. The sampler seed fixes the search sequence. Appendix~\ref{appendixA} (p.~\pageref{appendixA}) reports the selected configurations and notes that the Logistic Regression search sampled an inactive \texttt{l1\_ratio} while using the default $\ell_2$ penalty.

\subsection{FT-Transformer}

For FT-Transformer, five-fold cross-validation is replaced by a \textbf{single holdout validation split} (80/20 within DEV), because fitting a deep model for every fold would multiply the cost of each trial. This holdout-based choice is consistent with the evaluation design used in the FT-Transformer literature~\cite{gorishniy2021revisiting}, but it provides a different estimate of selection variability from the five-fold classical-model path.

Tuning also uses Optuna TPE with 50~trials and optimises PR-AUC on the validation split.
Each trial trains the model with early stopping (patience~=~15 epochs, monitoring validation PR-AUC).
The hyperparameter search space includes: number of Transformer blocks ($2$--$4$), token embedding dimension ($64$--$256$), number of attention heads ($4$--$8$), FFN dimension multiplier ($2/3$--$8/3$), attention and FFN dropout ($0.0$--$0.5$), learning rate ($10^{-5}$--$10^{-3}$), weight decay ($10^{-6}$--$10^{-3}$), and batch size ($256$--$1024$).
The decision threshold~$\tau$ is selected by maximising $F_2$ on the validation split (rather than the median across folds).

\subsection{General Notes}

For the BAF transfer study, tuning and fitting are performed on Base only. The saved Base-trained model, selected hyperparameters, and decision threshold are then applied unchanged to every Variant test distribution.

% ──────────────────────────────────────────────────────────────────────────
\section{Cross-Domain Generalisation}
\label{sec:cross_domain}
% ──────────────────────────────────────────────────────────────────────────

To assess robustness beyond the Base test distribution, the study performs transfer within the BAF suite. Variant~I changes group-size disparity, Variant~II changes group-wise prevalence disparity, Variant~III changes group-wise separability, Variant~IV changes prevalence disparity over time, and Variant~V changes separability over time. The target and \texttt{month} fields are excluded from modelling, as are the additional \texttt{x1} and \texttt{x2} columns in Variants~III and~V, so every model receives the same 30-feature input.

The protocol is as follows:
\begin{enumerate}
  \item \textbf{Base reference:} Train and evaluate each model on BAF Base using the standard holdout protocol.
  \item \textbf{Variant transfer:} Apply the saved Base-trained model to the holdout test set of each Variant~I--V \emph{without retraining or re-tuning}.
  \item \textbf{Threshold transfer:} Apply the decision threshold $\tau$ selected on Base validation unchanged to all Variant test sets.
  \item \textbf{Descriptive comparison:} Report PR-AUC, ROC-AUC, $F_1$, and $F_2$ on Base and each Variant. Differences are measured relative to the Base reference, not to separately trained in-domain Variant models.
\end{enumerate}

This is a controlled stress test of a fixed Base-trained model under the BAF suite's predefined conditions. It does not reproduce the full complexity of transfer between institutions or production concept drift. Figure~\ref{fig:cross_domain} illustrates the protocol.

\paragraph{Note on ULB \texorpdfstring{$\leftrightarrow$}{<->} BAF.}
Cross-dataset transfer between ULB~2013 and the BAF suite is not feasible due to incompatible feature spaces (PCA-anonymised vs.\ semantic features). The BAF internal transfer setting instead provides controlled Variant test distributions with a matched 30-feature modelling space.

% ═══════════════════════════════════════════════════════════════════════════
%  FIGURE 4.3 — Cross-Domain Generalisation Protocol (BAF Suite)
% ═══════════════════════════════════════════════════════════════════════════

\definecolor{clrBase}{HTML}{4472C4}      % Blue (BAF Base)
\definecolor{clrVariant}{HTML}{ED7D31}   % Orange (Variants)
\definecolor{clrDrop}{HTML}{CC0000}      % Red (degradation)

\par\medskip\noindent\begin{minipage}{\linewidth}\centering
\begin{tikzpicture}[
    >={Stealth[length=2.5mm]},
    node distance=0.6cm,
    basebox/.style={
        rectangle, rounded corners=3pt, minimum width=7cm, minimum height=1.1cm,
        fill=clrBase, text=white, font=\footnotesize\bfseries, align=center,
        drop shadow={shadow xshift=1pt, shadow yshift=-1pt}
    },
    varbox/.style={
        rectangle, rounded corners=3pt, minimum width=2.4cm, minimum height=1.1cm,
        text width=2.15cm,
        fill=clrVariant, text=white, font=\footnotesize\bfseries, align=center,
        drop shadow={shadow xshift=0.8pt, shadow yshift=-0.8pt}
    },
    compbox/.style={
        rectangle, rounded corners=3pt, minimum width=7cm, minimum height=0.8cm,
        fill=clrDrop!15, draw=clrDrop, line width=0.8pt,
        text=clrDrop, font=\footnotesize\bfseries, align=center
    },
    annot/.style={font=\scriptsize, text=black!60, align=center},
    arrow/.style={->, thick, black!60},
]

% ── BAF Base (source) ──
\node[basebox] (base)
    {BAF Base (Train)\\[-2pt]
     \fontsize{7}{8}\selectfont $N{=}800\text{K}$ DEV \;\;$\cdot$\;\; 30 features \;\;$\cdot$\;\; selected hyperparameters + $\tau$};

% ── Variants (targets) — single-line descriptions ──
\node[varbox, below=2.0cm of base, xshift=-5.2cm] (v1)
    {Variant I\\[-2pt] \fontsize{6.2}{7.2}\selectfont Group-size disparity};
\node[varbox, right=0.3cm of v1] (v2)
    {Variant II\\[-2pt] \fontsize{6.2}{7.2}\selectfont Prevalence disparity};
\node[varbox, right=0.3cm of v2] (v3)
    {Variant III\\[-2pt] \fontsize{6.2}{7.2}\selectfont Separability disparity};
\node[varbox, right=0.3cm of v3] (v4)
    {Variant IV\\[-2pt] \fontsize{6.2}{7.2}\selectfont Temporal prevalence};
\node[varbox, right=0.3cm of v4] (v5)
    {Variant V\\[-2pt] \fontsize{6.2}{7.2}\selectfont Temporal separability};

% ── Cross-domain label ──
\node[font=\scriptsize\bfseries, text=clrBase]
    at ($(base.south)!0.5!(v3.north)$) {Cross-domain transfer (no retraining)};

% ── Cross-domain arrows ──
\draw[arrow, clrBase] (base.south) -- ++(0,-0.5) -| (v1.north);
\draw[arrow, clrBase] (base.south) -- ++(0,-0.5) -| (v2.north);
\draw[arrow, clrBase] (base.south) -- ++(0,-0.5) -| (v3.north);
\draw[arrow, clrBase] (base.south) -- ++(0,-0.5) -| (v4.north);
\draw[arrow, clrBase] (base.south) -- ++(0,-0.5) -| (v5.north);

% ── Comparison ──
\node[compbox, below=1.0cm of v3] (comp)
    {Compare transferred metrics with the BAF Base reference};
\node[annot, below=0.2cm of comp] {PR-AUC \qquad ROC-AUC \qquad $F_1$ \qquad $F_2$};

\end{tikzpicture}
\captionof{figure}[Cross-domain generalisation protocol within the BAF suite]{Cross-domain generalisation protocol within the BAF suite. A model trained on BAF~Base is applied directly to Variants~I--V without retraining or threshold recalibration. The transferred metrics are compared descriptively with the Base reference; no separate in-domain Variant models are trained.}
\label{fig:cross_domain}
\end{minipage}\par\medskip

\clearpage

% ──────────────────────────────────────────────────────────────────────────
\section{Evaluation Metrics}
\label{sec:metrics}
% ──────────────────────────────────────────────────────────────────────────

The formal definitions of the metrics summarised below are given in Section~\ref{sec:tb:metrics}; this section states which metrics are reported and the specific operational choices made for their computation.

\paragraph{Primary and secondary summary metrics.}
The primary threshold-independent summary metric is average precision, computed with \texttt{average\_precision\_score} and reported as PR-AUC. ROC-AUC is a secondary ranking metric that supports comparison with prior work. The Brier score is reported for models that produce probability outputs as a summary of probabilistic accuracy. OCSVM produces an uncalibrated decision function rather than probabilities, so its clipped implementation value is neither reported nor interpreted as a Brier score.

\paragraph{Operating-point metrics.}
The primary operating-point metric is $F_2$, which assigns recall four times the weight of precision through the $\beta^2$ term and is the criterion used for threshold selection (Section~\ref{sec:threshold_study}). $F_1$ is reported alongside $F_2$ for the balanced-criterion view. Precision, recall, and the four confusion-matrix counts $\{TN, FP, FN, TP\}$ are reported in full to expose the underlying trade-offs.

\paragraph{Operational workload proxies.}
Two proxies for human review workload are reported: the alert rate $\frac{TP+FP}{N}$ (fraction of transactions flagged for review) and the ratio $\frac{FP}{TP}$ (false alerts per true fraud caught). These quantify how the chosen threshold rule translates into investigator effort. Precision@$k$ and Recall@$k$ for $k=0.5\%$ of scored transactions complement these by simulating an investigator who reviews only the highest-risk alerts under a fixed capacity budget.

\paragraph{Confidence intervals.}
Bootstrap $95\%$ confidence intervals ($B{=}1\,000$ iterations) are computed on the held-out test set for PR-AUC, ROC-AUC, and $F_2$. These intervals describe test-sample uncertainty conditional on each fitted model and split; they do not capture training-seed, split, or hyperparameter-search variability, and overlap between separate marginal intervals is not a formal paired significance test.

% ──────────────────────────────────────────────────────────────────────────
\section{Computational Cost and Reproducibility}
\label{sec:reproducibility}
% ──────────────────────────────────────────────────────────────────────────

All experiments are conducted on a single machine with an Intel Core i5-13600KF (14 cores), 16\,GB RAM, and an NVIDIA RTX~3070 (8\,GB VRAM). Classical ML models run on CPU via scikit-learn, LightGBM, and CatBoost; the FT-Transformer runs on GPU via PyTorch.

Each run logs its configuration (dataset identifier and SHA-256 hash, split seed, model parameters, preprocessing specification, imbalance strategy, and threshold rule), training and inference time, and hardware/software context (Python version and platform). Trained models and evaluation artefacts, including metric JSON files, PR-curve data, and confusion-matrix plots, are stored to support auditability and reproducibility. A common seed (42) is used for data splitting, Optuna sampling, and PyTorch initialisation; this controls the recorded runs but does not substitute for repeated-seed evaluation.

% ──────────────────────────────────────────────────────────────────────────
\section{Interpretability}
\label{sec:interpretability}
% ──────────────────────────────────────────────────────────────────────────

Interpretability is analysed through SHAP attributions~\cite{lundberg2017unified} on BAF~Base, whose features have semantic names. ULB~2013 features (PCA components V1--V28) are not interpreted in domain terms, so the substantive SHAP analysis is restricted to BAF. Section~\ref{sec:tb:shap} defines SHAP, LinearExplainer, TreeExplainer, and the Jaccard set-agreement measure.

The primary SHAP comparison covers LR, LGBM, and CatBoost. LinearExplainer is used for LR and TreeExplainer~\cite{lundberg2020local} for LGBM and CatBoost. For categorical features that are one-hot encoded during preprocessing, absolute SHAP values are aggregated back to the original feature level across one-hot columns. A GradientExplainer output was also generated for FT-Transformer using a 500-instance background and a 2,000-instance test sample. However, the explainer wrapper converts categorical inputs from floating-point tensors to integer embedding indices during its forward pass. This non-differentiable conversion prevents reliable gradient attribution to those categorical inputs. The FT-Transformer SHAP output is therefore retained only as an exploratory artefact and is excluded from the cross-model feature-consistency conclusions.

The analysis comprises four components. For LR, LGBM, and CatBoost, global feature importance is summarised as mean $|\text{SHAP}|$ across the full BAF Base test set, identifying the top-$k$ ($k=15$) features per model. Local explanations are produced for selected LGBM true-positive, false-positive, and false-negative cases. Cross-model consistency is quantified by pairwise Jaccard similarity between the top-10 feature \emph{sets} of LR, LGBM, and CatBoost; this statistic measures membership overlap, not rank order or attribution magnitude. Cross-variant stability applies the Base-trained LGBM model and TreeExplainer to each BAF test distribution and compares the resulting top-10 sets. It therefore characterises global set membership for one fixed model, not the validity of every local explanation under shift.

\paragraph{FT-Transformer attention diagnostic.}
Attention is analysed separately for the baseline BAF Base FT-Transformer. The diagnostic extracts the \texttt{[CLS]} query row from the final transformer layer, with PyTorch averaging attention weights across heads, and aggregates those weights over the 200,000 TEST observations. Normalised entropy, \texttt{[CLS]} self-attention, and mean categorical-versus-numerical attention describe how concentrated this specific last-layer average is. Head averaging may conceal head-specific structure, earlier layers are not included, and attention weights are not treated as faithful output attributions.

% ──────────────────────────────────────────────────────────────────────────
\section{Summary}
\label{sec:methodology_summary}
% ──────────────────────────────────────────────────────────────────────────

The methodology establishes a shared outer holdout, common reporting metrics, controlled imbalance comparisons, validation-only threshold selection, a fixed Base-to-Variant transfer protocol, and architecture-compatible SHAP conclusions. The classical models and FT-Transformer use different inner-validation and preprocessing paths, and these differences remain explicit. Guided by the mapping in Table~\ref{tab:rq_design_map}, Chapter~\ref{sec:chapterResults} (p.~\pageref{sec:chapterResults}) now presents the evidence in research-question order and interprets it within the documented boundaries.
